\documentclass[11pt,a4paper]{article}
\usepackage[utf8]{inputenc}
\usepackage{amsmath,amssymb,amsthm}
\usepackage{hyperref}
\usepackage{booktabs}
\usepackage{geometry}
\usepackage{lscape}
\geometry{margin=1in}

\newtheorem{theorem}{Theorem}[section]
\newtheorem{lemma}[theorem]{Lemma}
\newtheorem{proposition}[theorem]{Proposition}
\newtheorem{corollary}[theorem]{Corollary}
\newtheorem{definition}[theorem]{Definition}
\theoremstyle{remark}
\newtheorem{remark}[theorem]{Remark}

\title{The Benincasa--Dowker Action on Regular Lattices:\\
Exact Formulas, Variational Principles, and the Boundary of Dynamics}

\author{Thomas DiFiore}

\date{\today}

\begin{document}

\maketitle

\begin{abstract}
We present a comprehensive formalization, in Lean~4 with Mathlib, of the
combinatorial and variational properties of the Benincasa--Dowker (BD) action
on regular grid posets. The formalization comprises 58 files (approximately
8800 lines, zero \texttt{sorry}) and covers:
(I)~exact formulas including a total-variation decomposition,
discrete Gauss--Bonnet, and recursive dimensional structure;
(II)~variational principles including universal discrete concavity,
positive energy (flat space uniquely minimizes the BD action at fixed content),
and an ADM-type decomposition with proved sign structure;
(III)~continuum identities including an exact identification of the
2D renormalized action with total variation, a 3D weak-field identity relating
the spatial action to $\sum\delta_i^2$ under the content constraint, and
a spectral relation between the BD and Einstein--Hilbert quadratic forms
via the graph Laplacian;
(IV)~algebraic identities relevant to thermodynamic scaling, including the
parameter-free product $\alpha m^\alpha = m(\alpha m^{\alpha-1})$ underlying
the $TS = m/(d{-}2)$ relation, negative specific heat signatures, and
Bekenstein marginality at $d = 4$.
We also establish two negative results: the Euclidean BD overlap on regular grids
produces a non-vanishing second difference (precluding a massless graviton),
and the Lorentzian BD Hessian on regular grids has flat space as a saddle point.
These results delineate a precise boundary: regular lattices capture many
combinatorial and variational properties of the BD action but do not recover
the dynamical content (wave propagation, Einstein equations), which in the
causal-set program requires the Poisson-sprinkling framework of
Benincasa--Dowker~\cite{BD2010}.
\end{abstract}

\tableofcontents

%% ====================================================================
\section{Introduction}
%% ====================================================================

The Benincasa--Dowker (BD) action~\cite{BD2010} is a natural functional on
causal sets that, on Poisson sprinklings of a Lorentzian manifold, converges
in expectation to the Einstein--Hilbert action~\cite{BD2010,Dowker2013}.
In this paper we study the BD action on \emph{regular grid posets}, where
exact computation is possible.

Our results are organized into four parts (exact formulas, variational
principles, continuum identities, and algebraic scaling relations) plus
a section on negative results. All theorems are formalized in Lean~4
with the Mathlib library, with zero \texttt{sorry} statements.

\subsection*{Scope and limitations}

This work establishes combinatorial and variational properties of the BD action
on regular grids and their convex subsets. It does \emph{not} derive
gravitational dynamics (wave propagation, geodesic motion, Einstein equations)
from regular grids; as shown in Section~\ref{sec:negative}, this is not possible
in the regular-grid setting. The algebraic identities in Part~IV are structural
properties of the closed-form BD action; their connection to black-hole
thermodynamics requires physical interpretation that goes beyond what is
formalized here.

%% ====================================================================
\section{Definitions and Setup}\label{sec:defs}
%% ====================================================================

\begin{definition}[BD action]
For a finite poset $P$ with Hasse diagram $H(P)$, the \emph{Benincasa--Dowker
action} is
\begin{equation}\label{eq:bd}
  S_{\mathrm{BD}}(P) = |P| - |\mathrm{links}(H(P))|,
\end{equation}
where $|P|$ is the number of elements and $|\mathrm{links}|$ is the number of
covering relations (edges of the Hasse diagram).
\end{definition}

\begin{definition}[Grid poset]
The \emph{grid poset} $[m]^d = \{1,\ldots,m\}^d$ has the product partial order:
$(a_1,\ldots,a_d) \leq (b_1,\ldots,b_d)$ iff $a_i \leq b_i$ for all~$i$.
A \emph{convex subset} $S \subseteq [m]^d$ satisfies: if $a \leq c \leq b$
with $a,b \in S$, then $c \in S$.
\end{definition}

\begin{definition}[Variable-width grid]
A \emph{variable-width grid} with $T$ rows and widths $w_0, \ldots, w_{T-1}$
($w_i \geq 1$) is the convex subset of $[m] \times [\max w_i]$ consisting of
elements $(t, x)$ with $0 \leq x < w_t$, aligned at the left boundary.
The \emph{total variation} is $\mathrm{TV} = \sum_{i=0}^{T-2} |w_{i+1} - w_i|$.
\end{definition}

\begin{definition}[Renormalized action]\label{def:ren}
For a width profile $(w_1,\ldots,w_T)$ with \emph{content} $n = \sum w_i^{D-1}$
equal to that of a flat grid $T \cdot w^{D-1}$:
\begin{equation}
  S_{\mathrm{BD}}^{\mathrm{ren}} = S_{\mathrm{BD}}(\text{profile})
  - S_{\mathrm{BD}}(\text{flat at same content}).
\end{equation}
\end{definition}

\begin{definition}[Sorted profile]
A width profile $(w_1,\ldots,w_T)$ is \emph{sorted} if $w_1 \leq \cdots \leq w_T$.
Every profile can be sorted; the sorted version maximizes overlap and thus
minimizes $S_{\mathrm{BD}}$ among all permutations of the same multiset of widths
(Theorem~\ref{thm:valley}).
\end{definition}

\begin{definition}[Lorentzian BD action]
For a 2D grid with $T$ time slices and \emph{light-cone parameter} $c$, a
\emph{causal link} from $(t,x)$ to $(t{+}1,x')$ exists iff $|x - x'| \leq c$.
The \emph{Lorentzian BD action} replaces the all-to-all temporal overlap
$\min(w_i,w_{i+1})^2$ with the causal-link count within the light cone.
For $c < \min(w_i, w_{i+1})$, this count is $(2c{+}1)\min(w_i,w_{i+1}) - c(c{+}1)$,
which is \emph{linear} (not quadratic) in the smaller width.
\end{definition}

%% ====================================================================
\part{Exact Formulas}
%% ====================================================================

\section{The Exact 2D BD Formula}\label{sec:exact}

\begin{theorem}[Exact BD formula, \texttt{ExactBDFormula.lean}]
\label{thm:exact}
For a variable-width grid with $T$ rows and widths $w_0, \ldots, w_{T-1}$:
\begin{equation}
  2\,S_{\mathrm{BD}} = 2T - 2n + w_0 + w_{T-1} + \mathrm{TV},
\end{equation}
where $n = \sum w_i$.
\end{theorem}

\begin{proof}[Proof sketch]
The vertical links between adjacent rows satisfy
$2\,\mathrm{vLinks} + \mathrm{TV} = \mathrm{adjSum}$, where
$\mathrm{adjSum} = \sum(w_i + w_{i+1})$ telescopes to $2n - w_0 - w_{T-1}$.
The identity $\min(a,b) = (a + b - |a-b|)/2$ converts the overlap sum.
Full proof by \texttt{ring} after unfolding definitions.
\end{proof}

\begin{theorem}[Discrete Gauss--Bonnet, \texttt{DiscreteGaussBonnet.lean}]
For any convex subset~$S$:
$2\,S_{\mathrm{BD}}(S) = \sum_{x \in S} (2 - \deg(x))$,
where $\deg(x)$ counts covering relations involving~$x$.
\end{theorem}

\begin{proof}[Proof sketch]
By the handshaking lemma, $\sum\deg = 2|\mathrm{links}|$.
Then $2S_{\mathrm{BD}} = 2|S| - 2|\mathrm{links}| = \sum(2-\deg)$.
\end{proof}

\begin{theorem}[Euler characteristic, \texttt{EulerCharacteristic.lean}]
$S_{\mathrm{BD}}(S) = \chi_{\mathrm{graph}}(\mathrm{Hasse}(S)) = |V| - |E|$
(the graph Euler characteristic of the Hasse diagram).
\end{theorem}

\begin{theorem}[Recursive BD, \texttt{RecursiveBD.lean}]
\label{thm:recursive}
For the flat grid:
$S_{\mathrm{BD}}([m]^d) = m \cdot S_{\mathrm{BD}}([m]^{d-1}) - (m{-}1)m^{d-1}$.
Closed form: $S_{\mathrm{BD}}([m]^d) = (1{-}d)m^d + dm^{d-1}$.
Verified for $d = 2,3,4$ with $m = 3,4$.
\end{theorem}

\begin{proof}[Proof sketch]
Slice $[m]^d$ into $m$ copies of $[m]^{d-1}$ along the last coordinate.
Spatial links within slices contribute $m \cdot S_{\mathrm{BD}}([m]^{d-1})$.
Temporal links between adjacent slices contribute $(m{-}1) \cdot m^{d-1}$
(full overlap for equal slices). The closed form follows by induction on~$d$.
\end{proof}

\begin{theorem}[\texttt{ConvexIsVariableWidth.lean}]
Every convex subset of $[m] \times [n]$ is a variable-width grid.
$S_{\mathrm{BD}} = T' - \sum\mathrm{overlaps}$, and left-aligned profiles
minimize $S_{\mathrm{BD}}$ at fixed widths.
\end{theorem}

\section{Nested Total-Variation Structure}

\begin{theorem}[\texttt{NestedTV.lean}]
The BD action across dimensions has iterated TV structure:
$d = 2$ depends on $\mathrm{TV}^1$ (width TV),
$d = 3$ on $\mathrm{TV}^2$ (spatial-action TV),
$d = k$ on $\mathrm{TV}^{k-1}$.
Flat slices ($\mathrm{TV}^k = 0$) give minimum $S_{\mathrm{BD}}$ at each level.
\end{theorem}

\begin{theorem}[\texttt{NestedTVQuantitative.lean}]
For $d = 3$, $T = 3$: among all triples $(a,b,c)$ with $a^2+b^2+c^2 = n$
and $a,b,c \geq 1$, the minimum-TV triple minimizes $S_{\mathrm{BD}}$.
Verified exhaustively for $n \leq 49$ by \texttt{native\_decide}.
\end{theorem}

\begin{theorem}[\texttt{BDConjecture.lean}]
The discrete curvature $\sum(w_{i+1} - 2w_i + w_{i-1})$ telescopes to
boundary differences. In 2D, $S_{\mathrm{BD}}$ depends on TV, not on the
discrete Ricci scalar. This is consistent with 2D gravity being topological.
\end{theorem}

%% ====================================================================
\part{Variational Principles}
%% ====================================================================

\section{Universal Discrete Concavity}\label{sec:concavity}

Define $f_d(m) = (1{-}d)m^d + dm^{d-1}$ (the BD action on $[m]^d$).

\begin{theorem}[Universal concavity, \texttt{UniversalConcavityGeneral.lean}]
\label{thm:concavity}
For all $d \geq 2$ and all integers $n \geq 1$:
\begin{equation}
  f_d(n-1) + f_d(n+1) \leq 2\,f_d(n).
\end{equation}
\end{theorem}

\begin{proof}[Proof sketch]
Define $A_k(n) = (n{-}1)^k + (n{+}1)^k - 2n^k$ (the second finite difference of
$m^k$) and $B_k(n) = (n{+}1)^k - (n{-}1)^k$ (the first finite difference).
The recurrence $A_{k+1}(n) = n\,A_k(n) + B_k(n)$ holds by \texttt{pow\_succ}
and \texttt{ring}. The concavity defect decomposes as
\[
  D_d(n) = [d - (d{-}1)n] \cdot A_{d-1}(n) + (1{-}d) \cdot B_{d-1}(n).
\]
For $n \geq 2$: both terms are $\leq 0$ since $A_{d-1}(n) > 0$ (by induction),
$B_{d-1}(n) > 0$ (by \texttt{pow\_lt\_pow\_left\textsubscript{0}}), and
$d - (d{-}1)n \leq 2 - d \leq 0$. For $n = 1$: direct computation gives
$D_d(1) = (2{-}d)2^{d-1} - 2 \leq -2$.
The proof is 10 lines of Lean, using no induction on~$d$.
\end{proof}

Explicit defect polynomials for $d = 2, \ldots, 6$ are proved in
\texttt{GrandUnity.lean}; monotonicity of the defect in~$n$ is verified through
$d = 8$ in \texttt{UniversalConcavity.lean}.

\section{Positive Energy and Equal-Width Optimality}\label{sec:optimal}

\begin{theorem}[Sorted profile formula, \texttt{SortedProfileFormula.lean}]
\label{thm:sorted}
For $T$ sorted slices in dimension $D = k{+}2$ with content $n = \sum w_i^{D-1}$:
\begin{equation}
  S_{\mathrm{BD}} = -(D{-}1)\,n + (D{-}1)\sum_i w_i^{D-2} + w_T^{D-1}.
\end{equation}
At fixed content, $S_{\mathrm{BD}}$ is determined by $\sum w_i^{D-2}$ (the
``sum term'') and $w_T^{D-1}$ (the ``max-width penalty'').
\end{theorem}

\begin{proof}[Proof sketch]
Expand $S_{\mathrm{BD}} = \sum f_{D-1}(w_i) - \sum_{i<T} w_i^{D-1}$ (sorted
overlap). Substitute $f_{D-1}(w) = (2{-}D)w^{D-1} + (D{-}1)w^{D-2}$ and collect
terms. The result follows by \texttt{linarith}.
\end{proof}

\begin{theorem}[Equal-width optimality, \texttt{EqualWidthOptimal.lean}]
\label{thm:optimal}
For $D = 3$ and $T = 2, 3$: among sorted profiles with $\sum w_i^2 = Tw^2$ and
$w_i \geq 1$, the equal-width profile $(w,\ldots,w)$ uniquely minimizes
$S_{\mathrm{BD}}$.
\end{theorem}

\begin{proof}[Proof sketch]
From the sorted formula, the objective is $2\sum w_i + w_T^2$ at fixed
$\sum w_i^2$. For $T = 2$: the difference
$(w{-}a)(w{+}a{-}2) + 2(w_T{-}w) \geq 0$ since $a \leq w$ and
$w{+}a \geq 2$. For $T = 3$: case split on $b \leq w$ vs $b > w$, using
$(c{-}w)(c{+}w{+}2) \geq 2(2w{-}a{-}b)$. For $T \geq 3$, $w = 1$: the
content constraint forces all widths to~1.
\end{proof}

\begin{theorem}[Valley loss, \texttt{TVMinimization.lean}]
\label{thm:valley}
For sorted $a \leq b \leq c$: the ``valley'' arrangement $(b,a,c)$ has
strictly less overlap than the monotone arrangement $(a,b,c)$.
\end{theorem}

\begin{theorem}[Overlap bound, \texttt{ConvexLowerBound.lean}]
For any two intervals: $\mathrm{overlap}(\text{shifted}) \leq
\mathrm{overlap}(\text{left-aligned}) = \min(w_1, w_2)$.
\end{theorem}

\section{ADM-Type Decomposition}\label{sec:adm}

\begin{definition}
For a sorted profile with content $\sum w_i^2 = Tw^2$:
\begin{align}
  \Delta R_{\mathrm{spatial}} &= \sum w_i - Tw
  &&\text{(spatial curvature deficit, $\leq 0$ by Cauchy--Schwarz)},\\
  \Delta K_{\mathrm{extrinsic}} &= w_T^2 - w^2
  &&\text{(extrinsic curvature penalty, $\geq 0$ since $w_T \geq w$)}.
\end{align}
\end{definition}

\begin{theorem}[ADM structure, \texttt{ADMStructure.lean}]
\label{thm:adm}
For $D = 3$, $T = 2, 3$:
\begin{enumerate}
\item \emph{Exact decomposition}:
  $S_{\mathrm{BD}}^{\mathrm{ren}} = 2\,\Delta R + \Delta K$.
\item \emph{Sign structure}: $\Delta R \leq 0$ and $\Delta K \geq 0$.
\item \emph{ADM dominance}: $\Delta K \geq 2|\Delta R|$.
\item \emph{Rigidity} ($\iff$): $S_{\mathrm{BD}}^{\mathrm{ren}} = 0$
  if and only if $w_1 = \cdots = w_T = w$.
\end{enumerate}
\end{theorem}

\begin{proof}[Proof sketch]
(1) is algebraic (\texttt{linarith} from the sorted formula).
(2) follows from Cauchy--Schwarz and $w_T \geq w$.
(3) for $T = 2$: factor as $(w{-}a)(w{+}a{-}2) + 2(w_T{-}w)$ (both terms $\geq 0$).
(4) forward: if all equal, both terms vanish. Reverse: if the sum is zero and
both terms are $\geq 0$, each must be zero; $(w{-}a)(w{+}a{-}2) = 0$ with $a \geq 1$
forces $a = w$.
\end{proof}

%% ====================================================================
\part{Continuum Identities}
%% ====================================================================

\section{The 2D Bridge}\label{sec:2dbridge}

\begin{theorem}[2D bridge identity, \texttt{Bridge2DGeneral.lean}]
\label{thm:2dbridge}
Under fixed content ($\sum w_i = Tw$) and boundary ($w_0 = w_{T-1} = w$):
\begin{equation}
  2\,S_{\mathrm{BD}}^{\mathrm{ren}} = \mathrm{TV} = \sum_i |w_{i+1} - w_i|.
\end{equation}
Proved for $T = 2, 3, 4$ by the identity $|b{-}a| = a + b - 2\min(a,b)$.
\end{theorem}

\begin{theorem}[Displacement bound, \texttt{ContinuumLimitReal.lean}]
\label{thm:displacement}
For a $C^1$ function $w\colon [a,b] \to \mathbb{R}$ with $a \leq b$:
\begin{equation}
  |w(b) - w(a)| \leq \int_a^b |w'(t)|\,dt.
\end{equation}
Proved using Mathlib's \texttt{Convex.norm\_image\_sub\_le\_of\_norm\_deriv\_le}
and the Fundamental Theorem of Calculus (\texttt{integral\_deriv\_eq\_sub}).
\end{theorem}

\begin{remark}
Combining Theorems~\ref{thm:2dbridge} and~\ref{thm:displacement}:
for a Lipschitz profile $w(t)$ sampled at mesh $\ell$, each
$|w_{i+1}-w_i| \leq \int_{t_i}^{t_{i+1}} |w'|$, so
$S_{\mathrm{BD}}^{\mathrm{ren}} \leq \frac{1}{2}\int_0^L |w'(t)|\,dt$.
The convergence $S_{\mathrm{BD}}^{\mathrm{ren}} \to \frac{1}{2}\int|w'|$ as
$\ell \to 0$ for Lipschitz profiles follows from standard Riemann-sum theory
(uniform convergence of tagged sums for continuous integrands) but is
\emph{not formalized} in this Lean codebase. The formalization establishes
the algebraic identity (Theorem~\ref{thm:2dbridge}) and the one-step bound
(Theorem~\ref{thm:displacement}).
\end{remark}

\section{The 3D Weak-Field Structure}

\begin{theorem}[Separated-bump additivity, \texttt{SeparatedBumps.lean}]
\label{thm:bumps}
For an isolated upward bump of height $h$ at background width $w$ ($h \geq 0$):
\[
  \Delta S_{\mathrm{BD}} = f_2(w{+}h) - f_2(w) = -h(2w{-}2{+}h),
\]
where $f_2(w) = -w^2 + 2w$. The overlap at flat neighbors is unchanged:
$\min(w, w{+}h) = w$. For bumps separated by at least one flat slice,
the total $S_{\mathrm{BD}}^{\mathrm{ren}}$ equals the sum of individual costs.
Proved for 2 and 3 separated bumps (\texttt{two\_bumps\_additive},
\texttt{three\_bumps\_additive}).
\end{theorem}

\begin{theorem}[Weak-field identity, \texttt{WeakFieldLimit.lean}]
\label{thm:weakfield}
Under the content constraint $\sum(w{+}\delta_i)^2 = Tw^2$
(which implies $2w\sum\delta_i + \sum\delta_i^2 = 0$):
\begin{equation}
  w \cdot \sum_i \bigl(f_2(w{+}\delta_i) - f_2(w)\bigr) = -\sum_i \delta_i^2.
\end{equation}
The spatial part of $S_{\mathrm{BD}}^{\mathrm{ren}}$ is exactly $-(1/w)\sum\delta_i^2$.
Proved for $T = 2, 3$ by \texttt{linear\_combination -(w-1) * hc} where
\texttt{hc} is the content constraint.
\end{theorem}

\begin{theorem}[3D convergence structure, \texttt{ConvergenceTheorem.lean}]
\label{thm:3dconv}
For $C^1$ profiles $a(t)$ with $|a(t) - a_0| \leq M$: each Riemann-sum term
satisfies $(a(t_i) - a_0)^2 \leq M^2$, and the overlap oscillation satisfies
$|\min(a,b)^2 - a^2| \leq |b^2 - a^2| = |b-a| \cdot |b+a|$.
Informally, $w_0 \ell\, S_{\mathrm{BD}}^{\mathrm{ren}} \to -\int_0^L (a(t){-}a_0)^2\,dt$
as $\ell \to 0$. The algebraic bounds are formalized; the convergence statement
is not.
\end{theorem}

\section{Spectral Relation to Einstein--Hilbert}\label{sec:spectral}

\begin{theorem}[Spectral relation, \texttt{SpectralRelation.lean}]
\label{thm:spectral}
For periodic perturbations $u = (u_0, \ldots, u_{T-1})$, define the quadratic forms
\begin{align}
  F_{\mathrm{BD}}[u] &= \sum_i u_i^2 = \langle u, u \rangle,
  &&\text{(identity operator)} \\
  F_{\mathrm{EH}}[u] &= \sum_i (u_{i+1} - u_i)^2 = \langle u, ({-}\Delta_{\mathrm{graph}}) u \rangle,
  &&\text{(graph Laplacian)}
\end{align}
where $\Delta_{\mathrm{graph}} = 2I - 2S$ and $S$ is the cyclic shift.
The identity
\begin{equation}\label{eq:ibp}
  F_{\mathrm{EH}} = 2\,F_{\mathrm{BD}} - 2\sum_i u_i u_{i+1}
\end{equation}
is the discrete integration by parts. Proved for $T = 3, 4$ by \texttt{ring}.
\end{theorem}

Both quadratic forms share the eigenfunctions $e^{2\pi i k j / T}$;
$F_{\mathrm{BD}}$ has eigenvalue~1 for all~$k$, while $F_{\mathrm{EH}}$ has
eigenvalue $4\sin^2(\pi k/T) \to (2\pi k/L)^2$ in the continuum limit.

\begin{theorem}[Spectral bounds, \texttt{SpectralBDvsEH.lean}]
\label{thm:bounds}
\begin{enumerate}
\item $F_{\mathrm{EH}}[u] \leq 4\,F_{\mathrm{BD}}[u]$ for all $u$ (proved for $T \leq 4$).
\item $F_{\mathrm{BD}}[u] \leq C\,F_{\mathrm{EH}}[u]$ for mean-zero $u$ (Poincar\'e inequality,
  proved for $T = 2, 3$).
\end{enumerate}
\end{theorem}

%% ====================================================================
\part{Algebraic Scaling Relations}
%% ====================================================================

\section{Identities Relevant to Thermodynamic Scaling}\label{sec:thermo}

The closed-form action $f_d(m) = (1{-}d)m^d + dm^{d-1}$ gives rise to several
algebraic identities. Their connection to black-hole thermodynamics (in the sense
of Bekenstein--Hawking~\cite{Bekenstein1973,Hawking1975}) involves additional
physical interpretation; here we state only the algebraic content.

\begin{theorem}[\texttt{ParameterFreePrediction.lean}]
\label{thm:ts}
For any $\alpha \geq 1$ and $m \geq 1$:
$\alpha \cdot m^\alpha = m \cdot (\alpha \cdot m^{\alpha-1})$.
Interpreting $\alpha = d{-}2$ as the entropy exponent ($S \propto m^{d-2}$)
and $T = 1/((d{-}2) c m^{d-3})$ as the first-law temperature, this gives
$TS = m/(d{-}2)$, independent of the proportionality constant~$c$.
\end{theorem}

\begin{theorem}[\texttt{ThermodynamicSignatures.lean}]
For the BD action with $d \geq 4$: the exponent $d{-}3 > 0$ makes the
formal ``specific heat'' $C = -S/(d{-}3)$ negative, and the exponent
$d{-}4 = 0$ iff $d = 4$, corresponding to Bekenstein-type marginality.
\end{theorem}

\begin{theorem}[\texttt{RemnantMass.lean}]
$(d{-}1)(d{-}2) = 6$ at $d = 4$.
\end{theorem}

\begin{theorem}[\texttt{ContinuumScaling.lean}]
$S_{\mathrm{BD}} = (1{-}d) \cdot m^d + d \cdot m^{d-1}$: the BD action
is exactly linear in the volume $m^d$ and area $m^{d-1}$ terms.
For $m \geq 2$: $S_{\mathrm{BD}} \leq 0$ (the area term is subdominant).
\end{theorem}

\section{Spectral Gap and Partition Function}

\begin{theorem}[\texttt{SpectralGap.lean}]
The exact density of states for $[3]^2$:
$g({-}3)=1,\; g({-}2)=2,\; g({-}1)=9,\; g(0)=18,\;
g(1)=55,\; g(2)=27,\; g(3)=1$.
\end{theorem}

\begin{theorem}[\texttt{UniversalGap.lean}]
At $d = 2$: removing the minimum or maximum element from $[m]^2$ increases
$S_{\mathrm{BD}}$ by exactly~1 (spectral gap $\Delta = 1$).
Verified for $m = 2, 3, 4$.
\end{theorem}

\begin{theorem}[\texttt{SaddleDominance.lean}]
For the formal partition function $Z = \sum_{S} e^{-\beta S_{\mathrm{BD}}(S)}$
over convex subsets: at $\beta > m\ln 16$, the flat configuration dominates.
\end{theorem}

\begin{theorem}[\texttt{CylinderForced.lean}, \texttt{CylinderOptimal.lean}]
Convexity plus full boundary forces the full cylinder. The full cylinder
minimizes the 3D BD action among all convex subsets with those boundary conditions.
\end{theorem}

\section{Algebraic Structure}

\begin{theorem}[\texttt{ConvexityIFF.lean}]
$S \subseteq [m]^d$ is convex $\iff$ the restriction to~$S$ preserves
multiplication in the incidence algebra.
\end{theorem}

\begin{theorem}[\texttt{BulletproofRecovery.lean},
\texttt{RecoveryCounterexample.lean}]
The CSpec map (causal-spectrum functor) is injective under the $T_0$~axiom.
Without $T_0$: the V-poset $\{0,1,2\}$ with $0 \leq 2$ and $1 \leq 2$
gives a counterexample ($\uparrow 0 = \uparrow 1 = \{2\}$).
\end{theorem}

\begin{theorem}[\texttt{HolonomyComposition.lean}]
The transition-matrix composition $T_{p,q} \cdot T_{q,r} = e_q$ satisfies
12 composition laws, establishing holonomy as a functor.
\end{theorem}

\begin{theorem}[\texttt{DimensionalOrdering.lean},
\texttt{GammaOrdering.lean}]
$\gamma(\text{chain}) < \gamma(\text{grid}) < \gamma(\text{antichain})$ for
$m \geq 2$, where $\gamma = |CC|^{1/|P|}$ is the growth rate of convex subsets.
\end{theorem}

%% ====================================================================
\part{Negative Results}
%% ====================================================================

\section{The Boundary of Dynamics}\label{sec:negative}

\subsection{Euclidean grids: non-vanishing second difference}

\begin{theorem}[\texttt{LorentzianBD.lean}]
\label{thm:eucmass}
The Euclidean overlap $\min(a,b)^2$ has constant second difference:
\[
  (w{+}1)^2 + (w{-}1)^2 - 2w^2 = 2
\]
for all~$w$. By contrast, the Lorentzian overlap $(2c{+}1)\min(a,b) - c(c{+}1)$
has second difference $-(2c{+}1)$ (with respect to the smaller argument).
\end{theorem}

The non-vanishing second difference of the Euclidean overlap contributes
a constant diagonal term to the Hessian of $S_{\mathrm{BD}}$ at flat space.
In a continuum interpretation, this corresponds to a mass term for the
gravitational perturbation, precluding a massless graviton.
The experimental bound $m_g < 1.2 \times 10^{-22}$~eV~\cite{LIGO2017}
rules out a graviton mass at the lattice scale.

\subsection{Lorentzian grids: saddle instability}

Numerical computation of the Lorentzian BD Hessian at flat space gives
$H \approx -L_{\mathrm{path}}$ (negative graph Laplacian). The constrained
Hessian (projected onto the content-preserving subspace) has both positive
and negative eigenvalues, making flat space a \emph{saddle point} rather
than a minimum.

\subsection{Interpretation}

\begin{center}
\begin{tabular}{lcc}
\toprule
& \textbf{Combinatorial/variational} & \textbf{Dynamical} \\
\midrule
Regular grids (this paper) & Exact results & Does not recover EH \\
Poisson sprinklings~\cite{BD2010} & No exact results & Recovers EH \\
\bottomrule
\end{tabular}
\end{center}

Regular grids encode exact combinatorial and variational structure: counting,
entropy exponents, positive energy, and continuum identities.
The dynamical content of gravity (wave propagation, Einstein equations)
requires the Poisson-sprinkling framework~\cite{BD2010,Dowker2013,Surya2019},
where the random causal structure provides the degrees of freedom needed
for Lorentz-invariant dynamics.

%% ====================================================================
\section{Conclusion}
%% ====================================================================

We have presented a comprehensive formalization of the combinatorial and variational
properties of the Benincasa--Dowker action on regular grid posets. The formalization
is self-contained (58 Lean files, zero \texttt{sorry}) and covers exact formulas,
variational principles, continuum identities, algebraic scaling relations, and
negative results. The main findings are:

\begin{enumerate}
\item The BD action on grids has rich exact structure: total-variation decomposition,
  Gauss--Bonnet, Euler characteristic, recursive dimension formula.
\item The variational theory is complete in the grid setting: universal concavity,
  positive energy, ADM decomposition with sign structure and rigidity.
\item The 2D renormalized action equals half the total variation exactly;
  the 3D spatial action under the content constraint equals $-\sum\delta_i^2/w$;
  the BD and EH quadratic forms are spectrally related by the graph Laplacian.
\item Regular grids do not recover Einstein dynamics: the Euclidean overlap gives
  a mass term, and the Lorentzian Hessian has a saddle.
\end{enumerate}

The two-layer picture---regular grids for combinatorial structure, Poisson sprinklings
for dynamics---is the main conceptual conclusion.

%% ====================================================================
\begin{thebibliography}{99}

\bibitem{BD2010}
  D.~Benincasa and F.~Dowker,
  \emph{The scalar curvature of a causal set},
  Phys.\ Rev.\ Lett.\ \textbf{104}, 181301 (2010).

\bibitem{Dowker2013}
  F.~Dowker and L.~Glaser,
  \emph{Causal set d'Alembertians for various dimensions},
  Class.\ Quantum Grav.\ \textbf{30}, 195016 (2013).

\bibitem{Surya2019}
  S.~Surya,
  \emph{The causal set approach to quantum gravity},
  Living Rev.\ Relativ.\ \textbf{22}, 5 (2019).

\bibitem{Bekenstein1973}
  J.~D.~Bekenstein,
  \emph{Black holes and entropy},
  Phys.\ Rev.\ D \textbf{7}, 2333 (1973).

\bibitem{Hawking1975}
  S.~W.~Hawking,
  \emph{Particle creation by black holes},
  Commun.\ Math.\ Phys.\ \textbf{43}, 199 (1975).

\bibitem{LIGO2017}
  B.~P.~Abbott \emph{et al.}\ (LIGO/Virgo Collaboration),
  \emph{Tests of general relativity with GW150914},
  Phys.\ Rev.\ Lett.\ \textbf{116}, 221101 (2016).

\end{thebibliography}

\end{document}
